%% plantilla.tex - Informe de ejemplo para Simulación de Sistemas (ITBA)
%%
%% Estructura indicada por docs/GuiaInformes.md y AGENTS.md:
%%   Resumen; Introduccion; Modelo; Implementacion; Simulaciones;
%%   Resultados; Conclusiones; Referencias.
%%
%% Compilar desde este directorio:
%%   latexmk -pdf -interaction=nonstopmode -halt-on-error plantilla.tex
%%
%% Antes de entregar:
%%   - eliminar todos los \pendiente{...};
%%   - referenciar en el texto cada figura, tabla y ecuacion;
%%   - verificar ejes, unidades, barras de error y cifras significativas;
%%   - dejar en Referencias solamente las fuentes citadas.

\documentclass[11pt,a4paper]{article}
\usepackage{sdsreport}

% Datos de la caratula
\reporttitle{Trabajo Práctico Nro. X}
\reportsubtitle{Título específico del trabajo}
\reportgroup{Grupo 2}
\reportcommission{Comisión S2}
\reportmembers{%
  Arancibia Nicolás - Legajo pendiente\\
  Morales Agustín - Legajo pendiente\\
  Ronda Agustín - Legajo 64507}
\reportterm{Segundo cuatrimestre de 2026}
\reportdate{\today}
% El logo incluido en informes/ es el default. Para usar otro:
% \reportlogo{ruta/al/logo.png}

\begin{document}

\hypersetup{pageanchor=false}
\makereportcover
\setcounter{page}{1}
\hypersetup{pageanchor=true}

\begin{abstract}
\pendiente{redactar al final}. Resumir el problema, el modelo, el protocolo de
simulación y los resultados cuantitativos principales. El resumen no introduce
resultados que luego no aparezcan en el cuerpo del informe.
\end{abstract}

\section{Introducción}

Presentar el problema, el objetivo del estudio y las preguntas que se van a
responder. Citar los antecedentes necesarios, como se hace con la
Ref.~\cite{referencia-ejemplo}. Evitar adelantar resultados o describir en
detalle el sistema particular que se simuló.

\section{Modelo}

Exponer el modelo matemático general sin fijar la geometría, los parámetros de
una corrida ni el método numérico. Por ejemplo,
\begin{equation}
  \vect{r}_i(t+\Delta t)
  = \vect{r}_i(t) + \vect{v}_i(t)\,\Delta t,
  \label{eq:movimiento}
\end{equation}
donde \(\vect{r}_i\) es la posición de la partícula \(i\),
\(\vect{v}_i\) su velocidad, \(t\) el tiempo y \(\Delta t\) el paso temporal.
La Ec.~\eqref{eq:movimiento} también muestra la convención tipográfica:
vectores en negrita recta y escalares en itálica.

\section{Implementación}

Describir cómo se tradujo el modelo matemático al motor de simulación:
arquitectura, estructuras de datos, algoritmo de avance y criterio de
actualización. Incluir pseudocódigo o un diagrama cuando aclare relaciones que
serían difíciles de seguir en prosa. Esta sección no describe scripts de
graficación ni formatos de archivos de salida.

\section{Simulaciones}

\subsection{Sistema estudiado}

Definir el sistema particular: geometría, condiciones de contorno, parámetros
fijos, inputs barridos y outputs guardados. La Tabla~\ref{tab:parametros}
muestra un formato posible para registrar la configuración sin duplicar datos
que ya aparecen en figuras.

\begin{table}[htbp]
  \centering
  \caption{Configuración del estudio. Reemplazar cada marcador por el valor
  usado y agregar las unidades que correspondan.}
  \label{tab:parametros}
  \begin{tabular}{ll}
    \toprule
    Magnitud & Valor \\
    \midrule
    Geometría y contorno & \pendiente{descripción} \\
    Parámetros fijos & \pendiente{valores y unidades} \\
    Input barrido & \pendiente{rango y discretización} \\
    Realizaciones & \pendiente{cantidad de seeds distintas} \\
    Tiempo simulado & \pendiente{duración} \\
    \bottomrule
  \end{tabular}
\end{table}

\subsection{Observables y protocolo de medición}

Definir matemáticamente cada observable y explicar cómo se obtiene el escalar
que caracteriza una corrida. Si se promedia en estado estacionario, indicar
qué instantes entran en la suma y justificar el inicio del estacionario con
evoluciones temporales mostradas en Resultados.

Para \(R\) realizaciones independientes, un promedio y su desvío estándar
muestral pueden escribirse como
\begin{align}
  \langle \overline{A} \rangle
    &= \frac{1}{R}\sum_{k=1}^{R}\overline{A}_k,
    \label{eq:promedio-realizaciones}\\
  \sigma_A
    &= \sqrt{\frac{1}{R-1}\sum_{k=1}^{R}
       \left(\overline{A}_k-\langle\overline{A}\rangle\right)^2}.
    \label{eq:desvio-realizaciones}
\end{align}
Las barras de error deben indicar si representan \(\sigma_A\), el error
estándar u otra magnitud.

\section{Resultados}

Cada subsección debe sostener una secuencia de texto, figura y análisis. Para
cada input estudiado, conviene mostrar primero la dinámica característica,
luego la evolución temporal que valida el escalar y finalmente el observable
escalar en función del input.

\subsection{Dinámica y evolución temporal}

La Fig.~\ref{fig:evolucion} reserva el lugar para una evolución temporal. La
figura final debe marcar con una línea vertical el comienzo del estado
estacionario y mostrar el valor escalar que se usa en el estudio paramétrico.

\begin{figure}[htbp]
  \centering
  \fbox{%
    \parbox[c][5.0cm][c]{0.82\textwidth}{%
      \centering
      \pendiente{insertar una figura sin título interno, con ejes y unidades}}}
  \caption{Evolución temporal del observable. Explicar qué representa, cómo
  se identifica el estado estacionario y cuáles son los parámetros fijos.}
  \label{fig:evolucion}
\end{figure}

\subsection{Observable en función del input}

En la Fig.~\ref{fig:estudio-parametrico}, cada punto debe corresponder a un
promedio de realizaciones independientes y debe llevar su barra de error. Las
rectas que unen puntos se identifican como guías para el ojo; no se interpolan
los datos con una función sin fundamento teórico.

\begin{figure}[htbp]
  \centering
  \fbox{%
    \parbox[c][5.0cm][c]{0.82\textwidth}{%
      \centering
      \pendiente{insertar input vs. observable con puntos y barras de error}}}
  \caption{Observable estacionario en función del input. Informar cantidad de
  realizaciones, significado de las barras de error y parámetros fijos.}
  \label{fig:estudio-parametrico}
\end{figure}

\FloatBarrier

\section{Conclusiones}

Responder las preguntas planteadas en la Introducción con afirmaciones
cuantificadas y respaldadas por las figuras del informe. No incluir trabajos
futuros, resultados no mostrados ni explicaciones causales que no se hayan
puesto a prueba.

\begin{thebibliography}{9}

\bibitem{referencia-ejemplo}
Nombre Apellido,
``Título del trabajo'',
\textit{Nombre de la publicación}, vol. X, nro. Y, pp. 1--10 (2026).

\end{thebibliography}

\end{document}
